%=====================================================================
\chapter{Deployment, Maintenance and Performance}\label{ch:deploy}
%=====================================================================
\locator{5}{Part V --- Running and Extending the Shop}

\begin{outcomes}
By the end of this chapter you will be able to:
\begin{tight}
  \item \textbf{Measure} a shop's throughput and latency together, and
        \textbf{locate} its saturation point.
  \item \textbf{Explain} why throughput can fall as concurrency rises, and
        \textbf{say} what that means for a sale day.
  \item \textbf{Quantify} what the compiler cache is worth, and
        \textbf{order} optimisations by measured effect.
  \item \textbf{Distinguish} the cost of a backup from the cost of a restore,
        and \textbf{state} which one matters.
  \item \textbf{Assemble} a deployment checklist from this course's findings,
        with nothing on it that was not measured.
\end{tight}
\end{outcomes}

\begin{prereq}
\chref{stack}: images, containers, volumes and the Compose file.
\chref{security}: the hardening file and what it closed. \chref{gateway}:
reading a distribution rather than an average.
\end{prereq}

\begin{keyterms}
throughput $\bullet$ latency $\bullet$ concurrency $\bullet$ saturation point
$\bullet$ queueing $\bullet$ opcode cache $\bullet$ reverse proxy $\bullet$
backup $\bullet$ restore $\bullet$ recovery time $\bullet$ rollback $\bullet$
end of life $\bullet$ pinning $\bullet$ HSTS
\end{keyterms}

\section{Two Numbers, and Why One of Them Lies}

``How fast is the shop'' has two answers that move in opposite directions, and
quoting either alone is how capacity is mis-planned.

\begin{definitionbox}[Throughput, Latency, Saturation]
\term{Throughput} is requests served per second. It is a property of the
server, and it is the number in the sales brochure.

\smallskip
\term{Latency} is how long one request took. It is what the customer
experiences, and it is the number in the complaint.

\smallskip
The \term{saturation point} is the concurrency at which throughput stops
rising. Below it, extra customers are served from spare capacity. Above it,
extra customers join a queue --- so latency rises steeply while throughput
stays flat or \emph{falls}.
\end{definitionbox}

\begin{measured}[The same page, requested by more people at once]
From \texttt{code/m20\_deploy.py --load} against the stack in Appendix~A.
Sixty requests at each level, spread across that many concurrent clients.

\rowcolors{2}{white}{lightbg}
\begin{longtable}{@{}L{2.0cm}L{2.1cm}L{2.2cm}L{2.0cm}L{2.2cm}L{1.8cm}@{}}
\toprule
\rowcolor{primary!15}
\textbf{Clients} & \textbf{req/sec} & \textbf{Median} & \textbf{p95} &
\textbf{Worst} & \textbf{Failed} \\
\midrule
\endhead
1 & 11.0 & 81\,ms & 147\,ms & 158\,ms & 0 \\
\textbf{5} & \textbf{42.1} & 95\,ms & 181\,ms & 258\,ms & 0 \\
20 & \textbf{38.7} & \textbf{455\,ms} & 667\,ms & 814\,ms & 0 \\
\bottomrule
\end{longtable}

\smallskip
\textbf{Throughput peaked at five clients and then fell.} Twenty concurrent
customers were served at 38.7 requests per second --- \emph{fewer} than five
were. More customers, less work done.

\smallskip
\textbf{And the median response grew 5.6 times}, from 81\,ms to 455\,ms. Every
one of those twenty customers waited half a second for a page that one customer
alone gets in a twelfth of that.

\smallskip
\textbf{Nothing failed.} No error, no timeout, no log entry. The shop was past
its saturation point and the only observable was that it had become slow ---
which is the same observable as a badly written query, a slow network or a
customer's own machine. Load does not announce itself.

\smallskip
\textbf{The saturation point is the number to know before a sale day.} It is
five here, on a laptop, serving an uncached home page. What matters is not the
figure but that there \emph{is} one, that it is measurable in ten minutes, and
that past it every additional visitor makes the shop worse for everybody
already in it.
\end{measured}

\section{The Cheapest Optimisation Is Already a Setting}

Before any caching layer, any reverse proxy and any rewriting of queries, there
is one line of configuration worth measuring.

\begin{measured}[The same load, with the compiler cache off and on]
From \texttt{code/m20\_deploy.py --opcache}. PHP compiles source to bytecode;
the opcode cache keeps the result, so the files are compiled once rather than
on every request.

\rowcolors{2}{white}{lightbg}
\begin{longtable}{@{}L{2.0cm}L{2.8cm}L{2.8cm}L{2.2cm}L{3.4cm}@{}}
\toprule
\rowcolor{primary!15}
\textbf{Clients} & \textbf{req/sec off} & \textbf{req/sec on} &
\textbf{Gain} & \textbf{Median off $\to$ on} \\
\midrule
\endhead
1 & 8.1 & 10.1 & 1.3$\times$ & 109 $\to$ 82\,ms \\
\textbf{5} & \textbf{16.8} & \textbf{34.0} & \textbf{2.0$\times$} &
245 $\to$ 131\,ms \\
20 & 21.7 & 30.3 & 1.4$\times$ & 715 $\to$ 497\,ms \\
\bottomrule
\end{longtable}

\smallskip
\textbf{Twice the throughput at five clients, for one line of configuration.}
No code was changed, no caching layer added, no query rewritten.

\smallskip
\textbf{The gain is largest in the middle, and that is not a coincidence.} With
one client there is spare processor to waste, so the compilation is almost
free. At twenty, the queue dominates and the saving is diluted. At five --- the
saturation point --- the processor is the constraint, and removing work from it
is worth exactly what it costs.

\smallskip
\textbf{Check this before optimising anything else.} A shop running without it
is being benchmarked as a compiler rather than as a shop, and every conclusion
drawn from such a measurement is about the wrong program. \chref{stack}'s
Dockerfile installs it for this reason, and said so at the time.
\end{measured}

\begin{conceptbox}[The order to work in, from this course's own measurements]
\begin{tightnum}
  \item \textbf{Confirm the compiler cache is on.} 2$\times$, one line.
  \item \textbf{Minify the stylesheet} (\chref{ux}). 240\,KB off every page,
        one file, no visible change.
  \item \textbf{Index what the pages actually sort by} (\chref{catalogue}) ---
        and re-measure, because at 50{,}000 products the indexes made it
        \emph{slower}.
  \item \textbf{Cache whole pages} for visitors who are not logged in. The
        largest win available and the one that introduces staleness, so it
        comes after the three that do not.
  \item \textbf{Add servers.} Last, because everything above is cheaper, and
        because \chref{basket} showed that file-based sessions stop working the
        moment there are two.
\end{tightnum}
\end{conceptbox}

\section{Backups, and the Number Nobody Measures}

\begin{measured}[What a backup costs, and what a restore costs]
From \texttt{code/m20\_deploy.py --backup}.

\smallskip
\begin{tabular}{@{}L{5.0cm}L{3.2cm}L{5.4cm}@{}}
Database dump & \textbf{1.0\,s} & 572\,KB \\
\textbf{Database restore} & \textbf{6.4\,s} & \\
Shop files & 0.8\,s & 52\,MB in 4{,}024 files \\
\end{tabular}

\smallskip
\textbf{The restore takes six times as long as the backup}, and it is the only
one of the two numbers that matters --- because it is the one you will be
timing while the shop is down and somebody is standing behind you.

\smallskip
Those figures are for a demonstration catalogue. At \chref{catalogue}'s 50{,}000
products the dump was a different order of magnitude, and the restore grew with
it. \textbf{Measure your own, on your own data, and write the number down}: it
is your recovery time, and it is what you will be asked for.

\smallskip
\textbf{A backup nobody has restored is a file, not a backup.} The only way to
know a backup works is to restore it somewhere and look at the shop. Do it on a
schedule, not after an incident.
\end{measured}

\begin{alertbox}[Three things in the platform's own deployment advice that are
wrong]
Recorded because they are the first things a student will find, and all three
are dated rather than debatable.

\begin{tight}
  \item \textbf{Its own \texttt{Dockerfile} does not build}, because the one
        in the release contains the line
        \texttt{docker-php-ext-install zip \&\& \&\& docker-php-ext-enable
        zip} --- a shell syntax error. It has therefore never worked for
        anybody.
  \item \textbf{Its own compose file pins \texttt{mysql:5.7}}, which reached
        end of life in October 2023 and takes no further security fixes. This
        course's stack uses MariaDB 11 and \chref{stack} said why.
  \item \textbf{Its security guidance says to rename \texttt{.htaccess.txt}},
        which \chref{security} measured as doing nothing at all on a modern
        Apache.
\end{tight}

\smallskip
\textbf{Upstream artefacts are evidence, not authority.} Read them to learn what
the authors intended, then check whether they work.
\end{alertbox}

\section{Going Live}

\begin{examplebox}[The checklist, with a measurement behind every line]
Nothing here is advice this handout has not demonstrated.

\begin{tightnum}
  \item \textbf{Change all three default passwords} --- admin, database user,
        database root --- and raise the minimum password length from six
        (\chref{accounts}, \chref{security}).
  \item \textbf{Apply the hardening configuration} and re-run the battery. It
        moved 5 of 15 to 14 of 15 (\chref{security}).
  \item \textbf{Obtain a certificate, serve over HTTPS, and only then switch on
        HSTS.} Before the certificate it locks you out of your own shop.
  \item \textbf{Move \texttt{system/storage/} out of the web root}, rather than
        relying on a configuration file somebody will edit
        (\chref{security}).
  \item \textbf{Rename the admin directory} and the matching path in
        \texttt{admin/config.php}.
  \item \textbf{Confirm the compiler cache is on} --- 2$\times$, measured
        above.
  \item \textbf{Fix \texttt{config\_checkout\_payment\_address}}, without which
        no guest can complete a checkout (\chref{checkout}).
  \item \textbf{Populate \texttt{oc\_api\_ip}}, without which the admin panel
        cannot change an order (\chref{orders}).
  \item \textbf{Take a backup, restore it somewhere else, and time the
        restore.} Write the number down.
  \item \textbf{Know your saturation point} before the first sale day.
\end{tightnum}

\smallskip
\textbf{Four of those ten are faults in the shipped software} that every page
of the shop will hide from you. That is the argument of the whole course in one
list.
\end{examplebox}

\begin{pitfall}
\begin{tight}
  \item \textbf{Quoting throughput without latency.} At twenty clients this
        shop served 38.7 requests a second and every customer waited half a
        second.
  \item \textbf{Assuming more concurrency means more work done.} Past
        saturation it means less.
  \item \textbf{Benchmarking with the compiler cache off} and drawing
        conclusions about anything else.
  \item \textbf{Adding a caching layer before checking the free settings.}
  \item \textbf{Measuring the backup and never the restore.}
  \item \textbf{Never restoring a backup at all.}
  \item \textbf{Switching on HSTS before there is a certificate.}
  \item \textbf{Copying upstream's own Dockerfile or compose file.} One does not
        build and the other pins software that is out of support.
  \item \textbf{Deploying with the course's teaching passwords.}
\end{tight}
\end{pitfall}

\begin{lab}[Find where your shop falls over, then make it fall over later]
\begin{lstlisting}[language=bash]
# --- 1. Predict your saturation point -------------------------------
# How many concurrent customers before throughput stops rising?
# Write a number down before measuring.
cd code
python m20_deploy.py --load

# --- 2. Read both columns together ----------------------------------
# At the level where req/sec peaks, what is the median? And at
# twice that concurrency? Write one sentence a shop owner would
# understand.

# --- 3. Measure the free optimisation -------------------------------
python m20_deploy.py --opcache
# Explain why the gain is biggest in the MIDDLE row.

# --- 4. Apply Chapter 8's finding and re-measure --------------------
# Minify bootstrap.css, then run --load again. Did req/sec move?
# It should not, by much. Say why not, and what DID improve.

# --- 5. Time a restore ----------------------------------------------
python m20_deploy.py --backup

# --- 6. Now do it for real ------------------------------------------
# Seed 50,000 products (seed/catalogue.py), then time the dump and
# the restore again. Record both. That is your recovery time.
# Restore the demo catalogue afterwards.

# --- 7. Prove the backup works --------------------------------------
# docker compose down -v, up -d, restore your dump, and open the
# shop. If the products are not there, you did not have a backup.

# --- 8. Write the go-live checklist ---------------------------------
# Ten lines, each with the measurement or chapter that justifies it.
# Anything you cannot justify, take off the list.
\end{lstlisting}

\smallskip
\textbf{What to notice.} Step 7 is the only step that proves anything. Everyone
takes backups; the question that separates a professional from an amateur is
when they last restored one.
\end{lab}

\begin{checkpoint}
\begin{tightnum}
  \item At 5 concurrent clients the shop served 42.1 requests per second; at 20
        it served 38.7. Explain the mechanism, say what a shop owner should do
        with that information before a sale day, and give one change that would
        move the turning point.
  \item The compiler cache doubled throughput at 5 clients but gave only
        1.3$\times$ at 1 and 1.4$\times$ at 20. Explain both ends of that
        pattern, and say what it implies about when to benchmark.
  \item Your client's backup runs nightly and has never been restored. Write
        the case for a restore test in terms your client will act on, state
        exactly what you would do and how long it would take, and say what you
        would measure and record.
\end{tightnum}
\end{checkpoint}

\begin{chaptersummary}
\begin{tight}
  \item Throughput and latency are different questions and move in opposite
        directions. Quoting either alone mis-plans capacity.
  \item Measured: throughput peaked at 5 concurrent clients (42.1 req/sec) and
        \emph{fell} to 38.7 at 20, while the median response grew 5.6$\times$
        from 81\,ms to 455\,ms. Nothing failed and nothing was logged.
  \item The saturation point is measurable in ten minutes and is the number to
        know before a sale day.
  \item Measured: the opcode cache doubled throughput at the saturation point
        for one line of configuration, and the gain was largest in the middle
        --- spare capacity at one client, queueing at twenty.
  \item Optimise in measured order: compiler cache, then assets, then indexes
        (re-measuring, because at 50{,}000 products they hurt), then page
        caching, then servers.
  \item Measured: a 1.0\,s backup took 6.4\,s to restore. The restore is the
        number that matters, and a backup nobody has restored is a file.
  \item Upstream's own Dockerfile does not build, its compose file pins
        software out of support, and its security advice does nothing on a
        modern Apache. Artefacts are evidence, not authority.
  \item Four of the ten go-live checklist items are faults in the shipped
        software that every page of the shop hides from you.
\end{tight}
\end{chaptersummary}

\begin{reviewq}
  \item \textbf{(MCQ)} Going from 5 to 20 concurrent clients changed throughput
        from 42.1 to: (a)~168 req/sec; (b)~84 req/sec; (c)~38.7 req/sec;
        (d)~it failed.
  \item \textbf{(MCQ)} The opcode cache gave its largest gain at: (a)~1 client;
        (b)~5 clients; (c)~20 clients; (d)~the same at all three.
  \item \textbf{(Short)} Define the saturation point and explain why throughput
        can fall beyond it.
  \item \textbf{(Short)} Why is the restore time, rather than the backup time,
        the figure a business should be given?
  \item \textbf{(Short)} Explain why benchmarking a PHP application with the
        opcode cache disabled produces conclusions about the wrong program.
  \item \textbf{(Applied)} A shop expects 300 concurrent visitors for a
        two-hour sale. Its measured saturation point is 5 and its median
        response at 20 is 455\,ms. Set out what you would do, in order, with an
        expected effect for each step and a statement of which steps you could
        verify beforehand and which you could not.
  \item \textbf{(Applied)} Design the backup and recovery procedure for a shop
        taking 200 orders a day. State what is backed up, how often, where it
        is kept, how long you may lose, how long a restore takes, who performs
        it, how the procedure is tested, and how you would discover that a
        night's backup had silently failed.
\end{reviewq}
